\begin{env}[2.5.8.6]
\label{IV.2.5.8.6-fr}
\[
C'\xrightarrow{\sim}\operatorname{End}_{B'}(V'),\qquad
M'\xrightarrow{\sim}\operatorname{Hom}_{B'}(V',W')
\]
{\emergencystretch=2em
de las cuales la primera es un isomorfismo de $A'$-álgebras, y la segunda
forma con la primera un di-isomorfismo de $C'$-módulos a derecha.\par}
\end{env}

\begin{env}[2.5.8.7]
\label{IV.2.5.8.7-fr}
\emph{Reducción al caso en que $V=B_s$.} La hipótesis de que $V'$ y $W'$ son
$B'$-módulos isomorfos implica que $C'_d$ y $M'$ son $C'$-módulos a derecha
\emph{isomorfos} (2.5.6). Demostremos que, para probar (2.5.8), basta hallar un
elemento $u\in M$ que sea un \emph{generador} de $M$ como $C$-módulo a derecha. En
efecto, $u'=u\otimes1$ será un generador de $M'$ como $C'$-módulo a derecha;
ahora bien, en el caso (i), $V'$ y $W'$ son $A'$-módulos de tipo finito y, por
tanto, noetherianos, pues $A'$ es noetheriano; en el caso (ii), $V'$ y $W'$ son
$A'$-módulos (isomorfos) de presentación finita. Se puede, pues, aplicar (2.5.7)
a $A'$, $B'$, $V'$ y $W'$, y se deduce que $u'$ es un $B'$-isomorfismo de
$V'$ sobre $W'$. Pero, como $A'$ es fielmente plano sobre $A$, esto implica que
$u$ es biyectivo $(\mathbf{0}_{\mathrm I},6.4.1)$, es decir, la conclusión de (2.5.8).
Si se observa que, en el caso (i), $C$ y $M$ son $A$-módulos de tipo finito, y
que, en el caso (ii), $C$ es (como se ha visto en (2.5.8.1)) factor directo de
$V^n$ y, por tanto, un $A$-módulo proyectivo de tipo finito, se ve que la cuestión se reduce (cambiando las notaciones) a probar (2.5.8) en el caso particular en
que $V=B_s$, y que basta probar la existencia de un \emph{generador} del $B$-módulo
$W$. Se observará que entonces $B$ es un $A$-módulo \emph{de tipo finito}.
\end{env}

\begin{env}[2.5.8.8]
\label{IV.2.5.8.8-fr}
\emph{Reducción al caso en que $A$ y $A'$ son cuerpos y $B$ es una $A$-álgebra
simple de centro $A$.} Sea $\mathfrak r$ el radical del anillo semi-local $A$;
basta probar que $W/\mathfrak rW$ es un $(B/\mathfrak rB)$-módulo cíclico,
pues, si existe un homomorfismo sobreyectivo
$B_s/\mathfrak rB_s\to W/\mathfrak rW$, este da por composición un
homomorfismo $B_s\to B_s/\mathfrak rB_s\to W/\mathfrak rW$, el cual a su vez
(puesto que $B_s$ es un $B$-módulo libre) puede escribirse
$B_s\xrightarrow{f}W\to W/\mathfrak rW$, de modo que el homomorfismo
sobreyectivo considerado es
$f\otimes1:B_s\otimes_A(A/\mathfrak r)\longrightarrow W\otimes_A(A/\mathfrak r)$.
Como $W$ es un $A$-módulo de tipo finito, el lema de Nakayama demuestra que
$f$ es sobreyectivo (Bourbaki, \emph{Alg. comm.}, cap.~II, \S~3,
núm.~2, cor.~1 de la prop.~4). Si se pone $A_1=A/\mathfrak r$, $A'_1=A'\otimes_A A_1=A'/\mathfrak rA'$,
$B_1=B/\mathfrak rB=B\otimes_A A_1$, $W_1=W/\mathfrak rW=W\otimes_A A_1$,
las hipótesis (i) (resp. (ii)) siguen verificándose cuando se sustituye en ellas
$A$, $A'$, $B$, $V=B_s$, $W$ por $A_1$, $A'_1$, $B_1$,
$V_1=(B_1)_s$, $W_1$, respectivamente; además,
$V'_1=V'\otimes_A A_1=V_1\otimes_{A_1}A'_1$ y
$W'_1=W'\otimes_A A_1=W_1\otimes_{A_1}A'_1$ son $B'_1$-isomorfos (con
$B'_1=B'\otimes_A A_1=B_1\otimes_{A_1}A'_1$), y $A'_1$ es un $A_1$-módulo
fielmente plano. Se puede, pues, suponer, para demostrar (2.5.8), que $A$ es un
\emph{producto finito de cuerpos} (conmutativos). Como $B$ es un $A$-módulo de
tipo finito, es un anillo \emph{artiniano}; sea $\mathfrak R$ su radical. Bastará ahora
probar que $W/\mathfrak RW$ es un $(B/\mathfrak R)$-módulo cíclico, pues se ve,
como antes, utilizando el lema de Nakayama, que esto implica que $W$ es un
$B$-módulo cíclico; por otra parte, $W'/\mathfrak RW'$ es
$(B'/\mathfrak RB')$-isomorfo a $(B'/\mathfrak RB')_s$, y se tiene
$B'/\mathfrak RB'=(B\otimes_A A')\otimes_B(B/\mathfrak R)=(B/\mathfrak R)\otimes_A A'$
y, del mismo modo,
$W'/\mathfrak RW'=(W\otimes_A A')\otimes_B(B/\mathfrak R)=(W/\mathfrak RW)\otimes_A A'$.
Se puede, pues, suponer además que $\mathfrak R=0$, es decir, que $B$ es una
$A$-álgebra \emph{semisimple}.

Observemos ahora que, como $A$ es un producto finito de cuerpos $k_i$
$(1\leq i\leq n)$, $A'$ es producto directo de $k_i$-álgebras $A'_i$
$(1\leq i\leq n)$, y cada una de las $A'_i$ es anulada por los $k_j$ de índice
$j\ne i$; la hipótesis de que $A'$ es un $A$-módulo fielmente plano implica
que los $A'_i$ son $\ne0$; existe, por consiguiente, un cociente $A''_i$ de
$A'_i$ que es un cuerpo, y $A''$, producto directo de los $A''_i$, es un
$A$-módulo fielmente plano y un cociente de $A'$. Si se considera entonces el
anillo $B''=B'\otimes_{A'}A''=B\otimes_A A''$,
$W''=W'\otimes_{A'}A''=W\otimes_A A''$ es un $B''$-módulo isomorfo a
$B''_s$; es posible, pues, limitarse a demostrar (2.5.8) sustituyendo $A'$ por
$A''$, es decir, también se puede suponer que $A'$ es un \emph{producto finito
de cuerpos}.

Sea $Z$ el centro de $B$, que es un producto finito de cuerpos y un $A$-módulo
de tipo finito; observemos que $B$ y $W$ son $Z$-módulos de tipo finito,
proyectivos como todo $Z$-módulo; además, se tiene
$B'=B\otimes_Z Z'$ y $W'=W\otimes_Z Z'$, poniendo
$Z'=Z\otimes_A A'$, y $Z'$ es un $Z$-módulo fielmente plano. Se puede, pues,
sustituir $A$ por $Z$ en la hipótesis; en otras palabras, suponer que $A$ es el
\emph{centro} de $B$, siendo $B$ semisimple y $A'$ un producto finito de cuerpos. Si
los $k_i$ $(1\leq i\leq n)$ son los cuerpos componentes de $A$, $B$ es producto
directo de anillos simples $B_i$, siendo $k_i$ el centro de $B_i$, y $W$ es
suma directa de submódulos $W_i$ $(1\leq i\leq n)$, donde $W_i$ es anulado por
los $B_j$ de índice $j\ne i$; además, el razonamiento anterior muestra que
se puede suponer que $A'$ es un producto de cuerpos $k'_i$
$(1\leq i\leq n)$, donde $k'_i$ es una extensión de $k_i$ y es anulado por los
$k_j$ de índice $j\ne i$. La hipótesis de que $B'_s$ y $W'$ son $B'$-módulos
isomorfos implica entonces que, para
\oldpage[IV-2]{27}
todo $i$,
$(B_i\otimes_{k_i}k'_i)_s$ y $W_i\otimes_{k_i}k'_i$ son
$(B_i\otimes_{k_i}k'_i)$-módulos isomorfos; basta, pues, probar (2.5.8) cuando
$n=1$, es decir, en el caso en que $A$ y $A'$ son \emph{cuerpos} y $B$ es un álgebra
\emph{simple} de \emph{centro} $A$.
\end{env}

\begin{env}[2.5.8.9]
\label{IV.2.5.8.9-fr}
\emph{Fin de la demostración.} Se sabe entonces (Bourbaki, \emph{Alg.},
cap.~VIII, \S~5, prop.~6 y 8) que todo $B$-módulo es suma directa de módulos
isomorfos a un ideal minimal de $B$, y que dos $B$-módulos de rango finito
sobre $A$ son, por tanto, isomorfos si y solo si tienen el \emph{mismo rango} sobre
$A$. Por hipótesis, se tiene $[W':A']=[B'_s:A']$. Pero
$[W':A']=[W:A]$ y $[B'_s:A']=[B_s:A]$; por tanto, $[W:A]=[B_s:A]$, lo que
termina la demostración.
\end{env}
